\documentclass[10pt]{article}
\usepackage[T1]{fontenc}
\usepackage[utf8]{inputenc}
\usepackage[french]{babel}
\usepackage{amsmath,amssymb,mathrsfs,mathpazo}
\usepackage{geometry}
\geometry{a4paper,margin=2cm}
\usepackage{graphicx,array,tabularx,tabularray,colortbl,multirow,multicol,booktabs,enumitem,pifont,ulem,fancybox,diagbox,cancel}
\usepackage{pstricks,pstricks-add,pst-plot,pst-tree,pst-node,pst-text,pst-func,pst-math,pst-bspline,pst-3dplot,pst-eucl,pst-solides3d}
\let\pscalculate\undefined
\usepackage{pst-bezier}
\usepackage{tikz,pgfplots,tkz-tab}
\usetikzlibrary{arrows,arrows.meta,patterns,calc,positioning,trees,decorations.markings,shapes.geometric,shapes.misc,shapes.arrows,intersections,decorations.pathmorphing,decorations.pathreplacing,backgrounds,shadows,babel}
\pgfplotsset{compat=1.18}
\usepgfplotslibrary{fillbetween}
\usepackage[np]{numprint}
\usepackage{listings,makecell,pas-tableur,twoopt}
\let\barre\undefined
\usepackage[pstricks]{tablvar}
\definecolor{darkgreen}{RGB}{0,100,0}
\providecommand{\red}{\color{red}}
\providecommand{\blue}{\color{blue}}
\providecommand{\green}{\color{green}}
\providecommand{\black}{\color{black}}
\providecommand{\white}{\color{white}}
\providecommand{\gray}{\color{gray}}
\providecommand{\cyan}{\color{cyan}}
\providecommand{\magenta}{\color{magenta}}
\providecommand{\yellow}{\color{yellow}}
\providecommand{\violet}{\color{violet}}
\providecommand{\orange}{\color{orange}}
\providecommand{\euro}{\texteuro}
\providecommand{\EUR}{\texteuro}
\providecommand{\ieme}{\textsuperscript{e}}
\providecommand{\degres}{\ensuremath{^\circ}}
\providecommand{\degree}{\ensuremath{^\circ}}
\providecommand{\vv}[1]{\overrightarrow{#1}}
\providecommand{\Coord}[1]{#1}
\providecommand{\decosix}{}
\providecommand{\starredbullet}{\ensuremath{\bullet}}
\pagestyle{empty}
\setlength{\parindent}{0pt}
\usepackage{adjustbox}\DeclareUnicodeCharacter{2236}{:}\DeclareUnicodeCharacter{2190}{\ensuremath{\leftarrow}}
\providecommand{\R}{\mathbb{R}}
\providecommand{\N}{\mathbb{N}}
\providecommand{\Z}{\mathbb{Z}}
\providecommand{\Q}{\mathbb{Q}}
\providecommand{\D}{\mathbb{D}}
\providecommand{\C}{\mathbb{C}}
\providecommand{\e}{{\mathrm{e}}}
\providecommand{\vect}[1]{\overrightarrow{#1}}
\providecommand{\vectt}[1]{\overrightarrow{\mathrm{#1}}}
\providecommand{\vvt}[1]{\overrightarrow{\mathrm{#1}}}
\providecommand{\np}[2][]{{#2~#1}}
\providecommand{\vv}[1]{{\overrightarrow{#1}}}
\renewcommand{\d}{{\mathrm{d}}}
\providecommand{\degres}{{^\circ}}
\providecommand{\degre}{{^\circ}}
\providecommand{\blue}{\color{blue}}
\providecommand{\red}{\color{red}}
\providecommand{\enskip}{\quad}
\providecommand{\llbracket}{\lbrack\!\lbrack}
\providecommand{\rrbracket}{\rbrack\!\rbrack}
\providecommand{\euro}{{\text{€}}}
\providecommand{\no}{{\text{n°}}}
\providecommand{\interval}[2]{{[#1;#2]}}
\providecommand{\intervFO}[2]{{[#1;#2[}}
\providecommand{\dotfill}{\ldots}

\providecommand{\ds}{\displaystyle}
\providecommand{\textding}[1]{#1}
\providecommand{\Oij}{\(\mathrm{O};\vec{\imath},\vec{\jmath}\)}
\providecommand{\Oijk}{\(\mathrm{O};\vec{\imath},\vec{\jmath},\vec{k}\)}
\providecommand{\Ouv}{\(\mathrm{O};\vec{u},\vec{v}\)}
\providecommand{\intervFF}[2]{\ensuremath{[#1;#2]}}
\renewcommand{\euro}{\ifmmode\text{\texteuro}\else\texteuro\fi}
\DeclareUnicodeCharacter{00B0}{\ifmmode^{\circ}\else\textdegree\fi}
\DeclareUnicodeCharacter{20AC}{\ifmmode\text{\texteuro}\else\texteuro\fi}
% Énoncés : Éducation nationale. Transcription et corrigés : APMEP, Denis Vergès et contributeurs. Découpage et classement : CoopMaths, Sébastien Lozano, Éric Elter et contributeurs. Adaptation de la présentation : Maths Explorer.
% CC BY-SA : https://coopmaths.fr/www/about/
% https://coopmaths.fr/alea/static/bac/2021/tex/bac_2021_06_polynesie_2_cor.tex
\begin{document}
\noindent\textit{APMEP / CoopMaths. CC BY-SA. Présentation adaptée par Maths Explorer.}\par\medskip
\textbf{Commun à tous les candidats}

\medskip




















\begin{enumerate}
\item On construit un arbre pondéré modélisant la situation proposée:
\begin{center}
\pstree[treemode=R,nodesepB=3pt,levelsep=2.75cm]{\TR{}}
{\pstree{\TR{$M~~$}\taput{0,07}}
	{\TR{$T$} \taput{0,8}
	\TR{$\overline{T}$}\tbput{0,2} 
	}
\pstree{\TR{$\overline{M}~~$}\tbput{0,93}}
	{\TR{$T$} \taput{0,01}
	\TR{$\overline{T}$}\tbput{0,99} 
	}
}
\end{center}
\item 
	\begin{enumerate}
		\item 
On a $P(M \cap T) = P(M) \times P_M(T) = 0,07 \times 0,8 = 0,056$.
		\item 
		On a de même $P\left(\overline{M} \cap T\right) = P\left(\overline{M}\right) \times P_{\overline{M}}(T) = 0,93 \times 0,01 = \np{0,0093}$.
		
D'après la loi des probabilités totales :
		
$P(T) = P(M \cap T) + P\left(\overline{M} \cap T\right) = 0,056 + \np{0,0093} = \np{0,0653}$.
	\end{enumerate}
\item On calcule $P_T(M) = \dfrac{P(T \cap M)}{P(T)} = \dfrac{P(M \cap T)}{P(T)} = \dfrac{0,056}{\np{0,0653}} \approx \np{0,85758}$ soit $0,86$ à $10^{-2}$ près.






\item 


	\begin{enumerate}
		\item 
$X$ suit une loi binomiale de paramètres $n = 10$ avec $p = \np{0,0653}$.
		\item 


On a $P(X = 2) = \binom{10}{2}\times \np{0,0653}^2 \times (1 - \np{0,0653})^{10 - 2} = \binom{10}{2}\times \np{0,0653}^2 \times 0,9347^{8} =$

$ 45 \times \np{0,0653}^2 \times 0,9347^{8} \approx \np{0,1118}$, soit $0,11$ à $10^{-2}$ près.
	\end{enumerate}
\item 
On a $P(X \geqslant 1) = 1 - P(X = 0) = 1 - \binom{n}{0} \times \np{0,0653}^0 \times (1 - \np{0,0653})^{n} = 1 - \np{0,9347}^n$ et on veut que :

$P(X \geqslant 1) > 0,99 \iff 1 - \np{0,9347}^{n} > 0,99 \iff 0,01 > \np{0,9347}^{n}$ soit en prenant le logarithme népérien :

$\ln 0,01 > n\ln \np{0,9347} \iff \dfrac{\ln 0,01}{\ln \np{0,9347}}< n$.

Or $\dfrac{\ln 0,01}{\ln \np{0,9347}} \approx 68,2$. 

Il faut donc tester au moins 69 personnes au minimum.
\end{enumerate}

\end{document}
